% Generated draft: requires mathematical and source review.
\documentclass[11pt]{article}
\usepackage[margin=25mm]{geometry}
\usepackage{fontspec}
\setmainfont{Latin Modern Roman}
\usepackage{amsmath,amssymb,mathtools,unicode-math}
\setmathfont{Latin Modern Math}
\usepackage{enumitem,needspace,xurl,hyperref}
\setlength{\parindent}{0pt}
\setlength{\parskip}{5pt}
\setlength{\emergencystretch}{3em}
\newcommand{\N}{\mathbb{N}}
\newcommand{\Z}{\mathbb{Z}}
\newcommand{\Q}{\mathbb{Q}}
\newcommand{\R}{\mathbb{R}}
\newcommand{\C}{\mathbb{C}}
\newcommand{\F}{\mathbb{F}}
\newcommand{\A}{\mathbb{A}}
\newcommand{\PP}{\mathbb{P}}
\newcommand{\E}{\mathbb{E}}
\newcommand{\eps}{\varepsilon}
\newcommand{\dd}{\,\mathrm{d}}
\begin{document}
\section*{UnsolvedMath ID: 20000215}
\textbf{Source problem code:} catalogue entry
\subsection{Definitions and mathematical statement}
Eric Larson asks for possible fiber dimensions and vertical tangent-space dimensions of the evaluation map
\[
\pi:\mathcal M_{g,n}(\PP^r,\beta)\longrightarrow\operatorname{Conf}(\PP^r,n),
\qquad [f:C\to\PP^r,p_1,\ldots,p_n]\longmapsto(f(p_1),\ldots,f(p_n)).
\]
Here $g$ is the source genus, $\beta$ has projective degree $d$, and the target is the ordered configuration space of pairwise distinct image points; restrict to the locus where the images are distinct. The exact stable-map/open-locus and component conventions, as well as the ground field, are not fixed. The verification specifies $\PP^r$, resolving the extracted prose's $\PP^n$. The question especially concerns dominant evaluation maps. The source reports an explicit dominance characterization when $d\ge g+r$, without giving that characterization, and cites Stevens for canonical curves in that range. It asks for more than dominance.

On the smooth embedded locus, let $N_f$ be the normal bundle of $f$ and $D=\sum_i p_i$. The attached unreviewed aid reports
\[
T_{\mathrm{vert}}=H^0(C,N_f(-D)),\quad
\operatorname{coker}(d\pi)=\ker\bigl(H^1(C,N_f(-D))\to H^1(C,N_f)\bigr),
\]
\[
\dim T_{\mathrm{vert}}=(r+1)d-(r-3)(g-1)-(r-1)n+h^1(C,N_f(-D)).
\]
For a smooth plane curve of degree $e\ge3$, let $H$ be its hyperplane divisor, $N_e=e(e+3)/2$ and $g_e=(e-1)(e-2)/2$. The aid gives $\dim T_{\mathrm{vert}}=N_e-n+h^0(C,\mathcal O_C(D-3H))$. It claims no jumps for $n<3e$; at $n=3e$ the jump occurs exactly for $D\in|3H|$, equivalently when the marked divisor is a scheme-theoretic cubic intersection. Then it reports differential corank one and fiber dimension $g_e$, one above expected dimension $g_e-1$.

If $F_C$ defines $C$ and $Q$ is the cubic cutting out $D$, the reported linear system is
\[
\PP\bigl(kF_C+QH^0(\PP^2,\mathcal O(e-3))\bigr).
\]
To represent the smooth embedded fixed-configuration fiber, only its smooth members belong to the locus under discussion; the full projective system can contain singular curves. The aid calls it the whole smooth embedded fiber without explicitly making this open-locus restriction. Its plane result does not classify all stable-map fibers or tangent dimensions in every component.
\subsection{Short English statement}
Determine fiber and tangent dimensions when curves through marked projective points are evaluated at those points, especially in the dominant range. The aid gives normal-bundle formulas and a first plane-curve jumping locus at cubic intersections.
\subsection{Sources}
\begin{itemize}
\item Eric Larson, suggested Problem 1 and Remark 1.1 in AIM \emph{Degenerations in Algebraic Geometry}, September 7, 2016: \url{https://aimath.org/pastworkshops/degenalggeomproblems.pdf}. Canonical location: \texttt{aim-algebraic-geometry-notes.json}, notes[214]. The supplied verification reproduces the displayed evaluation map; this draft uses the extracted record, without consulting or compiling a PDF.
\item Original text reports Larson's recent dominance paper and work of Stevens for canonical curves, without titles or complete references. The separate literature field is empty.
\item The attached normal-bundle and cubic-jump aid is machine generated and not peer reviewed, with low novelty confidence; no full authoritative citations or proof artifacts are supplied.
\item Local registry: \texttt{data/unsolvedmath/problems.json}, permanent catalogue entry, exact source code \texttt{catalogue entry}.
\end{itemize}
\end{document}
